\chapter{Cadre d'administration}

\section{Objet du manuel}

Ce manuel décrit les opérations réservées aux deux niveaux d'administration de
ChargeTrackr :

\begin{itemize}[leftmargin=*]
  \item l'\rolebadge{Administrateur d'organisation}, responsable des employés,
  des stations, de la tarification et du suivi commercial de son organisation ;
  \item le \rolebadge{Super Admin}, responsable de l'ensemble de la plateforme,
  des organisations, des demandes de démonstration et des paramètres globaux.
\end{itemize}

Les fonctions destinées aux opérateurs, techniciens et clients sont détaillées
dans le \emph{Manuel utilisateur}. Les écrans et valeurs présentés dans ce
document correspondent à la version de référence indiquée sur la page de garde.

\section{Principes de gouvernance}

\subsection{Isolation des organisations}

Chaque station, connecteur, tarif, employé, alerte, intervention et opération
commerciale appartient à une organisation déterminée. Un administrateur ne
peut gérer que les ressources de sa propre organisation. Il ne peut ni
rechercher ni modifier les employés ou les actifs d'une autre organisation.

Le Super Admin dispose d'une vue transversale nécessaire à l'exploitation de la
plateforme. Cette capacité doit être utilisée uniquement pour les opérations
globales prévues et reste soumise au journal d'audit.

\securitynote{Ne jamais contourner une restriction d'organisation en utilisant
directement un identifiant aperçu dans une URL. Le serveur vérifie
l'autorisation et l'appartenance de la ressource pour chaque requête.}

\subsection{Séparation des responsabilités}

\begin{longtable}{|L{3.3cm}|L{5.4cm}|L{6.2cm}|}
  \hline
  \rowcolor{ctmint}
  \textbf{Rôle} & \textbf{Périmètre principal} & \textbf{Limites essentielles} \\ \hline
  \endfirsthead
  \hline
  \rowcolor{ctmint}
  \textbf{Rôle} & \textbf{Périmètre principal} & \textbf{Limites essentielles} \\ \hline
  \endhead
  Super Admin &
  Organisations, offres commerciales, demandes de démonstration, paramètres
  globaux, intégrations et audit &
  N'exécute pas les tâches quotidiennes propres à une organisation sauf besoin
  d'assistance contrôlé. \\ \hline
  Administrateur &
  Employés, clients liés, réseau de recharge, tarifs, rapports et abonnement de
  son organisation &
  Ne voit pas les ressources des autres organisations et ne modifie pas les
  paramètres globaux. \\ \hline
  Opérateur &
  Supervision opérationnelle, stations, OCPP, alertes, interventions et
  maintenance &
  Ne gère ni les comptes administrateurs ni les paramètres commerciaux
  globaux. \\ \hline
  Technicien &
  Consultation des stations, interventions assignées et comptes rendus &
  Ne crée pas de station et ne modifie pas les règles de gestion. \\ \hline
  Client &
  Recherche, recharge, paiement, reçus et adhésions &
  N'accède pas aux données internes d'exploitation. \\ \hline
\end{longtable}

\section{Accès et navigation}

L'authentification locale, l'activation par invitation, la connexion Google,
les notifications, la recherche globale, le profil et les paramètres
personnels suivent les mêmes principes que ceux décrits dans le
\emph{Manuel utilisateur}.

\begin{procedurebox}{Vérifier son périmètre avant une opération}
  \begin{enumerate}[leftmargin=*]
    \item Ouvrir le menu du profil et vérifier le rôle affiché.
    \item Pour un administrateur, vérifier également le nom de l'organisation.
    \item Contrôler le contexte de la page et les droits proposés par les
    boutons d'action.
    \item En cas d'incohérence, interrompre l'opération et contacter le support
    de la plateforme.
  \end{enumerate}
\end{procedurebox}

\section{Actions sensibles}

Les actions suivantes demandent une vérification particulière :

\begin{itemize}[leftmargin=*]
  \item invitation, désactivation ou suppression d'un compte ;
  \item création ou suppression d'une station ou d'un connecteur ;
  \item commande distante OCPP ;
  \item modification d'un tarif ou d'un plan ;
  \item suspension ou restauration d'une organisation ;
  \item modification d'un paramètre global.
\end{itemize}

\attention{Une donnée de démonstration ou un simulateur ne doit jamais être
présenté comme un équipement physique ou un paiement bancaire réel. Le statut
de l'intégration est visible dans l'espace Super Admin.}
